\documentclass[11pt, a4paper]{article}

\usepackage[T1]{fontenc}
\usepackage[utf8]{inputenc}
\usepackage[french]{babel}
\usepackage[margin=2cm]{geometry}
\usepackage{amsmath, amssymb}
\usepackage{booktabs}
\usepackage{enumitem}
\usepackage{titlesec}
\usepackage[expansion=false]{microtype}
\usepackage{xcolor}
\usepackage{graphicx}
\usepackage{float}
\usepackage{tikz}
\usepackage{hyperref}
\hypersetup{
    colorlinks=true,
    linkcolor=blue!60!black,
    urlcolor=blue!50!black,
    citecolor=blue!50!black,
    linktoc=all,
    bookmarks=true,
    bookmarksnumbered=true,
    pdftitle={IA Responsable, German Credit Dataset},
    pdfauthor={Gimenez, Fanchini, Cintra, Zhang, Hamadeh}
}
\usepackage{tcolorbox}
\tcbuselibrary{skins, breakable}

\definecolor{accent}{HTML}{2196F3}
\definecolor{fair}{HTML}{4CAF50}
\definecolor{alertc}{HTML}{d9534f}
\definecolor{neutral}{HTML}{FF9800}

\newtcolorbox{leconbox}[1][Leçon]{
    enhanced, breakable,
    colback=fair!8, colframe=fair!70!black,
    fonttitle=\bfseries\small, title={#1},
    left=4pt, right=4pt, top=3pt, bottom=3pt,
    boxrule=0.6pt
}
\newtcolorbox{attentionbox}[1][Attention]{
    enhanced, breakable,
    colback=alertc!8, colframe=alertc!70!black,
    fonttitle=\bfseries\small, title={#1},
    left=4pt, right=4pt, top=3pt, bottom=3pt,
    boxrule=0.6pt
}

\setlength{\parskip}{0.4em}
\setlength{\parindent}{0pt}
\titlespacing*{\section}{0pt}{1em}{0.4em}
\titlespacing*{\subsection}{0pt}{0.7em}{0.3em}
\setlist{nosep, leftmargin=1.5em}

\title{\textbf{IA Responsable, German Credit Dataset\\[0.3em]
\large Rapport accompagnant \texttt{responsiveAI-german-credit\_rendu.ipynb}}}
\author{%
IA708 --- Télécom Paris, 2026 \\[0.4em]
\normalsize \textbf{Étudiants :} \\
\normalsize Julien GIMENEZ \texttt{<julien.gimenez@telecom-paris.fr>} \\
\normalsize Hugo FANCHINI \texttt{<hugo.fanchini@telecom-paris.fr>} \\
\normalsize Paul CINTRA \texttt{<paul.cintra@telecom-paris.fr>} \\
\normalsize Yimou ZHANG \texttt{<yimou.zhang@telecom-paris.fr>} \\
\normalsize Zaher HAMADEH \texttt{<zaher.hamadeh@telecom-paris.fr>}
}
\date{}

\begin{document}
\maketitle
\vspace{-1em}

\begin{abstract}
\noindent Étude d'un modèle de scoring crédit (UCI German Credit, 1000 lignes) sous l'angle de l'IA responsable. On évalue une régression logistique baseline sur trois axes : équité (DP, EO, DI, audit intersectionnel genre $\times$ âge), interprétabilité (SHAP global et local, détection de proxies univariés et multivariés), et quantification d'incertitude (bootstrap $B = 500$ sur le baseline, $B = 100$ sur trois variantes ; en substitution à la robustesse adversariale, conformément à la consigne reçue). La démarche compare systématiquement trois versions du modèle : V1 avec attributs sensibles, V2 sans (baseline), V3 sans + proxies univariés ($|r| > 0.15$) retirés. Stack : \texttt{scikit-learn}, \texttt{fairlearn}, \texttt{shap} ; exécution reproductible via \texttt{papermill}. Résultat principal : V3 améliore simultanément la performance (AUC 0.785 $\to$ 0.790) et l'équité par attribut isolé (DI âge 0.69 $\to$ 0.87, au-dessus du seuil légal EEOC ; $|\Delta\text{DP}_{\text{age}}|$ de 0.125 à 0.044) sans intervention in/post-modèle, mais ne résout pas le biais intersectionnel (DI = 0.48 sur genre $\times$ âge).
\end{abstract}

\section*{Utilisation de l'IA dans ce projet}

Dans le cadre de ce projet, nous avons utilisé ponctuellement des outils d'intelligence artificielle générative, notamment des assistants conversationnels de type ChatGPT ou Claude. Leur utilisation n'a pas eu vocation à remplacer le travail d'analyse, de conception ou d'expérimentation réalisé par le groupe, mais plutôt à servir de support méthodologique et technique à certaines étapes du projet.

Ces outils ont principalement été employés comme moteurs de recherche augmentés, afin d'obtenir rapidement des explications générales sur certaines notions du cours, sur les méthodes de machine learning équitable et interprétable, ou encore sur des points techniques liés à leur mise en \oe uvre. Les informations obtenues par ce biais ont systématiquement été confrontées aux supports de cours, aux documents fournis dans le cadre du projet, ainsi qu'aux sources scientifiques ou documentaires pertinentes lorsque cela était nécessaire. L'objectif était donc de faciliter la compréhension initiale, sans considérer les réponses générées comme des références définitives.

Les outils d'IA ont également été utilisés comme aide ponctuelle à la programmation. Ils ont pu contribuer à générer certaines portions de code simples, à proposer des structures de fonctions, à expliquer des erreurs rencontrées lors de l'exécution, ou encore à suggérer des pistes de correction pour des bugs. Cependant, le code final intégré au projet a été relu, adapté, testé et validé par les membres du groupe. Les choix méthodologiques, les expérimentations réalisées, l'interprétation des résultats et les conclusions du rapport relèvent du travail du groupe.

Enfin, ces outils ont servi d'appui pour la structuration et la rédaction du rapport. À partir du contenu, des résultats et des orientations définis au préalable par le groupe, ils ont pu aider à reformuler certaines explications, à améliorer la clarté de certains passages ou à générer une première version structurée de certaines sections. La version finale du rapport a toutefois été relue, corrigée et complétée par le groupe afin de garantir sa cohérence avec le travail effectivement réalisé.

\tableofcontents
\vspace{1em}

% ════════════════════════════════════════════════════════════════
\section{Données et préparation}

\subsection{Dataset et attributs sensibles}

Le dataset UCI Statlog German Credit (Hofmann, 1994) comporte 1000 lignes et 20 attributs (6 numériques, 14 catégoriels), avec une cible binaire (30 \% de défauts, 70 \% de bons payeurs). Les codes catégoriels \texttt{Axx} sont traduits en libellés français lors du chargement, afin que les noms de features apparaissant dans les figures SHAP et dans les audits soient interprétables.

Deux attributs sensibles sont retenus pour l'audit : le genre, dérivé de \texttt{personal\_status\_sex} (690 hommes, 310 femmes) ; et l'âge binarisé au seuil de 25 ans, suivant la convention Kamiran \& Calders (2012) reprise par AIF360 (Bellamy et al.\ 2018), qui aboutit à 810 \emph{older} et 190 \emph{younger}. Le choix de ce seuil est justifié dans la littérature par l'idée que les jeunes adultes n'ont pas encore eu l'opportunité de bâtir un historique de crédit ; cette particularité statistique ne doit pas se traduire mécaniquement par une exclusion algorithmique.

Le RGPD (article 9) et la directive européenne 2000/43/CE classent l'âge et le sexe comme caractéristiques protégées : un modèle de scoring crédit ne peut s'en servir pour discriminer, qu'il s'agisse d'une utilisation directe ou indirecte via des proxies.

\subsection{Biais bruts}

\begin{table}[H]
\centering
\begin{tabular}{lcc}
\toprule
\textbf{Groupe} & \textbf{P(bon)} & \textbf{P(défaut)} \\
\midrule
Homme    & 0.72 & 0.28 \\
Femme    & 0.65 & 0.35 \\
\midrule
Older    & 0.73 & 0.27 \\
Younger  & 0.58 & \textbf{0.42} \\
\bottomrule
\end{tabular}
\caption{Taux de défaut conditionnel par groupe sensible (données brutes).}
\end{table}

Les taux de défaut diffèrent fortement entre groupes. Un modèle qui apprend à reproduire la distribution conditionnelle observée encodera mécaniquement ces écarts. C'est la justification statistique de l'audit fairness qui suit.

\subsection{Détection de proxies}
\label{sec:proxies-early}

Un proxy est une feature corrélée à un attribut sensible, suffisamment informative pour permettre au modèle de reconstruire indirectement cette information protégée. La conséquence pratique est qu'un simple retrait des attributs sensibles ne suffit pas à supprimer le biais : tant que les proxies sont en place, le modèle continue de capter le signal discriminatoire.

On utilise deux mesures complémentaires. Le premier est la corrélation univariée $|r|$ de Pearson entre chaque feature encodée et l'attribut sensible binaire, avec un seuil empirique conservateur $|r| > 0.15$ (la littérature retient des valeurs comprises entre 0.1 et 0.2). Le second est une régression logistique multivariée features $\to$ attribut sensible : son AUC mesure la \emph{découvrabilité} de l'attribut depuis le reste des features, et une valeur $> 0.65$ indique un signal proxy multivarié significatif (au-dessus du hasard à 0.5).

\begin{table}[H]
\centering
\begin{tabular}{lcc}
\toprule
\textbf{Attribut sensible} & \textbf{AUC multivariée} & \textbf{Nb features identifiées comme proxy} ($|r|>0.15$) \\
\midrule
Genre & 0.69 & 8 \\
Âge   & \textbf{0.86} & 15 \\
\midrule
Union & --- & 17 \\
\bottomrule
\end{tabular}
\caption{Détection de proxies : régression multivariée et corrélation univariée.}
\end{table}

Les proxies identifiés en univarié sont, pour l'âge, \texttt{housing\_locataire} ($|r|=0.33$), \texttt{present\_employment\_since\_$\geq 7$ans} (0.27), \texttt{credit\_history} (0.24), \texttt{job\_management} (0.24) ; et pour le genre, \texttt{number\_of\_people\_liable} (0.23), \texttt{present\_employment\_since} (0.21), \texttt{housing\_locataire} (0.21). L'AUC multivariée features $\to$ âge atteint 0.86, ce qui confirme qu'un modèle pourrait \emph{reconstruire} l'âge à partir du reste des variables avec une bonne fiabilité, et que retirer seulement \texttt{age\_in\_years} laisserait intact l'essentiel du signal discriminatoire.

\begin{leconbox}[Conséquence opérationnelle]
On constitue une troisième variante du jeu d'entraînement (\texttt{X\_tr\_np}, \texttt{X\_va\_np}, \texttt{X\_te\_np}), notée V3, où les 17 features ainsi identifiées sont retirées en plus des attributs sensibles. Cette variante est comparée tout au long du rapport à V1 (avec attributs sensibles) et V2 (sans, baseline), notamment dans la section~\ref{sec:perf-3}.
\end{leconbox}

\subsection{Préparation}

Le pipeline de prétraitement utilise un \texttt{ColumnTransformer} (Pedregosa et al.\ 2011) qui applique un \texttt{StandardScaler} aux variables numériques et un \texttt{OneHotEncoder} aux variables catégorielles. On obtient 56 features finales après encodage, ramenées à 39 dans la variante V3. Les données sont séparées par un split stratifié 60 / 20 / 20 (\texttt{random\_state = 42}) en train, validation et test, garantissant la reproductibilité des résultats.

% ════════════════════════════════════════════════════════════════
\section{Modèle baseline}
\label{sec:baseline}

\subsection{Choix et entraînement}

Le modèle baseline est une régression logistique L2 (\texttt{sklearn.linear\_model.LogisticRegression}, solver \texttt{liblinear}, $C = 1$), choisie pour sa simplicité, l'expressivité de ses coefficients (chaque feature a une contribution linéaire signée) et la disponibilité d'un Linear SHAP exact. L'entraînement converge en 6 itérations effectives sur le train (600 lignes). Le seuil de classification $\tau = 0.18$ est calibré sur le set de validation pour maximiser la balanced accuracy, ce qui équilibre les erreurs sur la classe minoritaire (les défauts).

\subsection{Performance globale}

\begin{table}[H]
\centering
\begin{tabular}{lc}
\toprule
\textbf{Métrique} & \textbf{Test} \\
\midrule
AUC ROC                          & 0.785 \\
Average Precision                & 0.607 \\
Balanced Accuracy                & 0.669 \\
\midrule
Confusion (TN / FP / FN / TP)    & 66 / 74 / 8 / 52 \\
TPR (rappel défaut)              & 86.7\,\% \\
PPV (précision)                  & 41.3\,\% \\
\bottomrule
\end{tabular}
\end{table}

\begin{figure}[H]
\centering
\includegraphics[width=0.95\linewidth]{../outputs/10_baseline_perf.png}
\caption{Matrice de confusion (gauche), courbe ROC AUC = 0.785 (centre), courbe Precision-Recall AP = 0.607 (droite).}
\end{figure}

\subsection{Coût métier (matrice UCI)}

La page UCI documente explicitement une matrice de coûts asymétriques pour ce dataset : \og{} It is worse to class a customer as good when they are bad (5), than it is to class a customer as bad when they are good (1) \fg{} (source : \url{https://archive.ics.uci.edu/dataset/144/statlog+german+credit+data}). Cette asymétrie reflète une intuition bancaire : laisser passer un mauvais payeur (FN) entraîne une perte financière supérieure au manque à gagner d'un bon client refusé (FP).

Cette information est utilisée pour comparer deux choix de seuil. Le premier, $\tau_{\text{balacc}} = 0.18$, maximise la balanced accuracy et conduit à un coût test de 114 unités. Le second, $\tau_{\text{cost}} = 0.14$, minimise directement $5 \cdot \text{FN} + 1 \cdot \text{FP}$ et donne un coût test de 122. Contre-intuitivement, c'est le seuil balanced accuracy qui se montre le plus économe sur cet échantillon, parce que sa TPR (taux de détection des défauts) est identique mais sa précision sur les non-défauts un peu meilleure. On retient $\tau = 0.18$ pour la suite du rapport.

% ════════════════════════════════════════════════════════════════
\section{Audit d'équité}

\subsection{Métriques utilisées}

Les calculs s'appuient sur \texttt{fairlearn.metrics.MetricFrame} (Bird et al.\ 2020), qui calcule chaque métrique par sous-groupe puis prend l'écart maximum. Trois quantités résument l'équité du modèle. Le \emph{Demographic Parity gap} $|\Delta\text{DP}|$ mesure l'écart entre taux de sélection (taux d'octroi du crédit) entre groupes ; il est dans $[0, 1]$, et plus petit est mieux. L'\emph{Equalized Odds gap} $|\Delta\text{EO}|$ (Hardt et al.\ 2016) mesure l'écart maximum entre TPR ou FPR par groupe ; il garantit que les vrais bons payeurs ont la même chance d'être acceptés indépendamment du groupe. Le \emph{Disparate Impact ratio} (DI) est le rapport entre le taux d'octroi du groupe le moins favorisé et celui du groupe le plus favorisé ; la règle US des 80~\% (EEOC, 1978) considère qu'un DI inférieur à 0.8 fait présumer une discrimination indirecte.

\begin{attentionbox}[{Théorème d'impossibilité (Chouldechova 2017 ; Kleinberg et al.\ 2017)}]
Lorsque les taux de défaut différent entre groupes (ce qui est le cas ici, 27 \% vs 42 \% entre \emph{older} et \emph{younger}), les critères DP et EO ne peuvent être satisfaits simultanément par un même modèle. Toute intervention qui force DP à zéro fait remonter EO, et inversement. Il faut donc choisir explicitement quel critère est prioritaire, en fonction du contexte légal et opérationnel.
\end{attentionbox}

\subsection{Audit par axe (baseline sans attributs sensibles)}

\begin{table}[H]
\centering
\begin{tabular}{lccc}
\toprule
\textbf{Attribut} & $|\Delta\text{DP}|$ & $|\Delta\text{EO}|$ & \textbf{DI} \\
\midrule
Genre & 0.010 & 0.058 & 0.97 \\
Âge   & \textbf{0.125} & 0.066 & \textbf{0.69} \\
\bottomrule
\end{tabular}
\end{table}

\begin{leconbox}[Constat]
Sur le genre, biais quasi nul après retrait de \texttt{personal\_status\_sex}. Sur l'âge, $|\Delta\text{DP}| = 0.125$ persiste et DI = 0.69 viole la règle US 80\,\%. Point d'alerte principal pour un audit de conformité.
\end{leconbox}

\subsubsection*{Audit sur les 3 variantes (V1, V2, V3)}

On compare V1 = avec attributs sensibles, V2 = sans (baseline), V3 = sans + proxies univariés ($|r| > 0.15$) retirés.

\begin{table}[H]
\centering
\small
\begin{tabular}{llccc}
\toprule
\textbf{Variante} & \textbf{Attribut} & $|\Delta\text{DP}|$ & $|\Delta\text{EO}|$ & \textbf{DI} \\
\midrule
V1 (avec sensibles) & Genre & 0.081 & 0.100 & 0.78 \\
V1 (avec sensibles) & Âge   & 0.138 & 0.115 & 0.62 \\
\midrule
V2 (sans sensibles) & Genre & 0.010 & 0.058 & 0.97 \\
V2 (sans sensibles) & Âge   & 0.125 & 0.066 & 0.69 \\
\midrule
\textbf{V3 (sans + proxies)} & Genre & 0.050 & 0.044 & 0.86 \\
\textbf{V3 (sans + proxies)} & Âge   & \textbf{0.044} & \textbf{0.045} & \textbf{0.87} \\
\bottomrule
\end{tabular}
\caption{Audit d'équité sur les 3 variantes du modèle, par attribut.}
\end{table}

\begin{figure}[H]
\centering
\includegraphics[width=0.95\linewidth]{../outputs/17_fairness_triple.png}
\caption{Comparaison visuelle des 3 variantes par attribut sensible (gauche : genre, droite : âge). Pointillé rouge : seuil EEOC 80\,\%.}
\end{figure}

\begin{leconbox}[Sur l'âge, V3 fait passer DI au-dessus du seuil EEOC]
Retirer les proxies divise par 3 le biais $|\Delta\text{DP}_{\text{age}}|$ (0.125 vers 0.044) et fait passer DI au-dessus du seuil légal pour la première fois (0.69 vers 0.87).
\end{leconbox}

\subsection{Audit intersectionnel}

\begin{table}[H]
\centering
\begin{tabular}{lccc}
\toprule
\textbf{Sous-groupe} & \textbf{n} & \textbf{Approval rate} & \textbf{TPR} \\
\midrule
female\_older   & 36  & 0.44 & 0.85 \\
male\_older     & 120 & 0.38 & 0.85 \\
female\_younger & 25  & 0.28 & 0.92 \\
male\_younger   & 19  & \textbf{0.26} & 0.86 \\
\bottomrule
\end{tabular}
\end{table}

\begin{table}[H]
\centering
\begin{tabular}{lccc}
\toprule
\textbf{Axe} & $|\Delta\text{DP}|$ & $|\Delta\text{EO}|$ & \textbf{DI} \\
\midrule
Genre seul     & 0.010 & 0.058 & 0.97 \\
Âge seul       & 0.125 & 0.066 & 0.69 \\
\textbf{Intersectionnel} & \textbf{0.181} & \textbf{0.222} & \textbf{0.59} \\
\bottomrule
\end{tabular}
\end{table}

\begin{figure}[H]
\centering
\includegraphics[width=0.95\linewidth]{../outputs/14_intersectional.png}
\end{figure}

\begin{leconbox}[Surprise]
Les \emph{male\_younger} (n = 19) ont le taux d'octroi le plus bas (0.26), pas les \emph{female\_younger} comme l'intuition le suggérait. L'intersection multiplie le biais par 1.5 par rapport à l'analyse par axe seul.
\end{leconbox}

\subsubsection*{Intersectionnel sur les 3 variantes}

\begin{table}[H]
\centering
\begin{tabular}{lccc}
\toprule
\textbf{Variante} & $|\Delta\text{DP}|$ & $|\Delta\text{EO}|$ & \textbf{DI} \\
\midrule
V1 (avec sensibles)  & 0.207 & 0.222 & 0.44 \\
V2 (sans sensibles)  & 0.181 & 0.222 & 0.59 \\
V3 (sans + proxies)  & 0.221 & 0.243 & 0.48 \\
\bottomrule
\end{tabular}
\caption{Audit intersectionnel (genre $\times$ âge) sur les 3 variantes.}
\end{table}

\begin{attentionbox}[{L'équité par attribut n'implique pas l'équité intersectionnelle}]
Contrairement aux axes isolés, retirer les proxies n'améliore pas l'intersectionnel : V3 garde un DI = 0.48 (légèrement pire que V2). Les proxies retirés étaient utiles pour les sous-groupes minoritaires (jeunes femmes notamment). \textbf{L'équité intersectionnelle ne résulte pas mécaniquement de l'équité par attribut isolé.}
\end{attentionbox}

% ════════════════════════════════════════════════════════════════
\section{Atténuation des biais}

\subsection{Pré-traitement : reweighing (Kamiran \& Calders 2012)}

L'idée de Kamiran \& Calders (2012) est de corriger le déséquilibre observé sur le train en pondérant chaque exemple par
$$w_i = \frac{P(S = s_i) \cdot P(Y = y_i)}{P(S = s_i, Y = y_i)},$$
ce qui revient à imposer l'indépendance statistique $S \perp Y$ \emph{a posteriori} sur les marges. La méthode est \emph{model-agnostic} : on ré-entraîne le même \texttt{LogisticRegression} en passant \texttt{sample\_weight=w}. Sur ce dataset, l'effet est limité (V2 : $|\Delta\text{DP}|$ âge passe de 0.125 à 0.132 ; le reweighing ne modifie pas significativement les frontières apprises), ce qui s'explique par le fait que la dépendance résiduelle $S \perp Y$ est faible une fois \texttt{personal\_status\_sex} retiré, et que les vrais leviers du biais sont les proxies multivariés que le reweighing ne touche pas.

\subsection{Post-traitement : seuils par groupe (Hardt et al.\ 2016)}

L'approche post-traitement est plus directe : au lieu d'un seuil global $\tau$, on calibre un seuil $\tau_g$ par groupe sensible sur la validation, en l'ajustant pour égaliser le taux de sélection. Sur le baseline V2, cette procédure aboutit à $\tau_{\text{older}} = 0.16$ et $\tau_{\text{younger}} = 0.30$. Ce résultat peut sembler paradoxal puisqu'il \emph{augmente} le seuil pour les jeunes, mais il s'explique exactement par le déséquilibre observé : comme le taux de défaut de base est plus élevé chez les jeunes (42 \% vs 27 \%), à seuil global constant ils sont plus souvent refusés. Pour égaliser le taux d'octroi global, il faut accepter à score plus élevé chez les jeunes, ce qui revient à n'admettre que les meilleurs dossiers de ce groupe.

\subsection{Comparaison des 4 configurations}

Le tableau ci-dessous met en regard les 4 configurations (baseline seul, reweighing seul, baseline + post-processing, reweighing + post-processing) sur le baseline V2. Les chiffres sont produits par la cellule \texttt{compare-cell} du notebook.

\begin{table}[H]
\centering
\small
\begin{tabular}{llcccc}
\toprule
\textbf{Attr.} & \textbf{Configuration} & \textbf{AUC} & \textbf{BalAcc} & $|\Delta\text{DP}|$ & $|\Delta\text{EO}|$ \\
\midrule
Genre & Baseline        & 0.785 & 0.669 & 0.010 & 0.058 \\
Genre & Reweighing      & 0.785 & 0.665 & 0.006 & 0.050 \\
Genre & Baseline + PP   & 0.785 & 0.669 & 0.105 & 0.196 \\
Genre & Reweighing + PP & 0.785 & 0.669 & 0.128 & 0.230 \\
\midrule
Âge   & Baseline        & 0.785 & 0.669 & 0.125 & 0.066 \\
Âge   & Reweighing      & 0.784 & 0.640 & 0.132 & 0.098 \\
Âge   & Baseline + PP   & 0.785 & 0.644 & \textbf{0.060} & \textbf{0.135} \\
Âge   & Reweighing + PP & 0.784 & 0.650 & 0.069 & 0.221 \\
\bottomrule
\end{tabular}
\end{table}

\begin{figure}[H]
\centering
\includegraphics[width=0.95\linewidth]{../outputs/02_tradeoff.png}
\end{figure}

\begin{attentionbox}[{Théorème d'impossibilité \emph{en action}}]
Sur l'âge, le post-processing réduit $|\Delta\text{DP}|$ de \textbf{0.125 à 0.060} mais \textbf{double} $|\Delta\text{EO}|$ (0.066 vers 0.135). Sur le genre où le biais initial est négligeable, le post-processing introduit un biais artificiel : le seuil par groupe calibré sur 200 lignes de validation ne généralise pas.
\end{attentionbox}

\subsubsection*{Comparaison 4 configurations $\times$ 3 variantes (cible : âge)}

\begin{table}[H]
\centering
\small
\begin{tabular}{llccccc}
\toprule
\textbf{Variante} & \textbf{Configuration} & \textbf{AUC} & \textbf{BalAcc} & $|\Delta\text{DP}|$ & $|\Delta\text{EO}|$ & \textbf{DI} \\
\midrule
V1 & Baseline        & 0.781 & 0.656 & 0.138 & 0.115 & 0.62 \\
V1 & Reweighing      & 0.777 & 0.674 & 0.112 & 0.057 & 0.71 \\
V1 & Baseline + PP   & 0.781 & 0.633 & \textbf{0.027} & 0.121 & \textbf{0.92} \\
V1 & Reweigh + PP    & 0.777 & 0.651 & 0.053 & 0.178 & 0.86 \\
\midrule
V2 & Baseline        & 0.785 & 0.669 & 0.125 & 0.066 & 0.69 \\
V2 & Reweighing      & 0.784 & 0.640 & 0.132 & 0.098 & 0.63 \\
V2 & Baseline + PP   & 0.785 & 0.644 & \textbf{0.060} & 0.135 & 0.86 \\
V2 & Reweigh + PP    & 0.784 & 0.650 & 0.069 & 0.221 & 0.81 \\
\midrule
\textbf{V3} & \textbf{Baseline}        & \textbf{0.790} & 0.652 & \textbf{0.044} & \textbf{0.045} & \textbf{0.87} \\
V3 & Reweighing      & 0.788 & 0.658 & 0.047 & 0.025 & 0.87 \\
V3 & Baseline + PP   & 0.790 & 0.608 & 0.124 & 0.250 & 0.71 \\
V3 & Reweigh + PP    & 0.788 & 0.604 & 0.140 & 0.225 & 0.69 \\
\bottomrule
\end{tabular}
\caption{Toutes les configurations $\times$ les 3 variantes, sur l'attribut âge.}
\end{table}

\begin{leconbox}[V3 baseline est le meilleur état global]
Sur V1 et V2, le PP corrige $|\Delta\text{DP}|$ au prix d'une hausse de $|\Delta\text{EO}|$ (théorème d'impossibilité). Sur V3, le PP est \emph{contre-productif} : il déborde de l'équilibre déjà atteint et fait remonter $|\Delta\text{DP}|$ de 0.044 à 0.124. Le meilleur compromis global reste \textbf{V3 baseline} sans intervention.
\end{leconbox}

% ════════════════════════════════════════════════════════════════
\section{Interprétabilité}

\subsection{SHAP global}

Pour la régression logistique, les valeurs de Shapley admettent une expression exacte (Lundberg \& Lee 2017, Corollary 1) :
$$\phi_i(x) = w_i \cdot (x_i - \mathbb{E}[x_i]),$$
où $w_i$ est le coefficient appris et $\mathbb{E}[x_i]$ la moyenne d'arrière-plan estimée sur le train. Cette forme close (\texttt{shap.LinearExplainer} avec \texttt{feature\_perturbation = "interventional"}) évite l'échantillonnage Monte Carlo et donne une décomposition exacte du logit prédit, additive sur les features et conservative ($f(x) = \phi_0 + \sum_i \phi_i$).

\begin{figure}[H]
\centering
\includegraphics[width=0.55\linewidth]{../outputs/03_shap.png}
\end{figure}

Les quatre features dominantes en importance moyenne ($\overline{|\phi_i|}$) sont \texttt{checking\_status} (0.67), \texttt{credit\_history} (0.49), \texttt{installment\_rate} (0.39) et \texttt{savings\_account\_bonds} (0.36). Toutes sont des variables financières légitimes, et leur ordre est confirmé par la permutation importance (figure \texttt{04\_shap\_vs\_perm.png} du notebook), qui mesure la perte d'AUC lorsqu'on permute chaque feature.

\subsection{SHAP local : explication individuelle}

Le RGPD article 22 reconnaît aux personnes concernées un droit à une explication concernant les décisions automatisées qui les affectent. Le SHAP local répond à cette obligation en décomposant la prédiction individuelle en contributions additives par feature, ce qui permet de produire une justification fidèle au modèle, et non une rationalisation \emph{post hoc}.

\begin{figure}[H]
\centering
\includegraphics[width=0.95\linewidth]{../outputs/09_shap_local.png}
\caption{SHAP local pour deux clients : haut risque à gauche (id 832), bas risque à droite (id 159). Barres rouges = contributions vers défaut, barres vertes = vers non-défaut.}
\end{figure}

Pour le client à haut risque (id 832, score 0.95, vrai défaut), la prédiction est dominée par trois facteurs : \texttt{credit\_amount} = 14\,555 ($\phi$ = +1.18), \texttt{checking\_status} négatif ($\phi$ = +0.81) et \texttt{duration\_in\_month} = 45 mois ($\phi$ = +0.66). La somme de ces trois contributions explique l'écart entre le logit prédit et la valeur moyenne $\phi_0$. Pour le client à bas risque (id 159, score 0.01), toutes les contributions individuelles tirent vers le non-défaut (compte sain, historique impeccable, taux d'échéance bas, épargne disponible), et la décision est confortée. Cette lecture est directement utilisable dans une procédure de motivation des décisions algorithmiques.

\subsection{Comparaison SHAP avant vs après retrait}

\begin{figure}[H]
\centering
\includegraphics[width=0.95\linewidth]{../outputs/08_shap_compare.png}
\end{figure}

Lorsqu'on compare le modèle V1 (avec attributs sensibles) au modèle V2 (sans), les features qui absorbent l'effet retiré sont identifiables par l'augmentation de leur SHAP : \texttt{present\_employment\_since} ($+0.036$), \texttt{telephone} ($+0.022$), \texttt{credit\_history} ($+0.020$). Cette redistribution est la signature observable de l'effet proxy, et confirme graphiquement le diagnostic quantitatif (AUC features $\to$ âge $=$ 0.86).

\subsubsection*{SHAP sur les 3 variantes}

\begin{figure}[H]
\centering
\includegraphics[width=0.95\linewidth]{../outputs/18_shap_triple.png}
\caption{Importance SHAP moyenne (top 12 features), comparée sur V1, V2, V3.}
\end{figure}

Les trois variantes pointent les mêmes features dominantes (\texttt{checking\_status}, \texttt{credit\_history}, \texttt{credit\_amount}), ce qui indique que le signal prédictif principal reste financier dans tous les cas. La redistribution s'observe sur les features de second rang : en passant de V1 à V3, l'importance migre depuis les proxies retirés (\texttt{housing}, \texttt{present\_employment\_since}, \texttt{job}) vers d'autres variables financières restées dans le modèle (\texttt{savings\_account\_bonds}, \texttt{installment\_rate}). En d'autres termes, V3 s'appuie davantage sur des features réellement explicatives et moins sur des variables potentiellement discriminatoires.

\subsection{Performance et équité : avec / sans / sans + proxies}
\label{sec:perf-3}

Le tableau ci-dessous récapitule, pour chacune des trois variantes, la performance globale et l'équité sur l'attribut âge (le seul des deux à présenter un biais significatif).

\begin{table}[H]
\centering
\begin{tabular}{lccccc}
\toprule
\textbf{Modèle} & \textbf{AUC} & \textbf{BalAcc} & $|\Delta\text{DP}_{\text{age}}|$ & $|\Delta\text{EO}_{\text{age}}|$ & \textbf{DI}$_{\text{age}}$ \\
\midrule
Avec attributs sensibles            & 0.781 & 0.656 & 0.138 & 0.115 & 0.62 \\
Sans attributs sensibles (baseline) & 0.785 & 0.669 & 0.125 & 0.066 & 0.69 \\
\textbf{Sans sensibles + proxies}   & \textbf{0.790} & 0.652 & \textbf{0.044} & \textbf{0.045} & \textbf{0.87} \\
\bottomrule
\end{tabular}
\caption{Triple comparaison : performance et équité selon le retrait progressif des features.}
\end{table}

\begin{figure}[H]
\centering
\includegraphics[width=0.85\linewidth]{../outputs/16_metrics_triple.png}
\caption{Performance (gauche, plus haut = mieux) et équité âge (droite, plus bas = mieux) selon les trois variantes.}
\end{figure}

\begin{leconbox}[Résultat marquant]
Retirer également les proxies identifiés \emph{améliore} l'AUC ($+0.005$) et fait passer DI âge de 0.69 à \textbf{0.87}, au-dessus du seuil légal EEOC ; $|\Delta\text{DP}_{\text{age}}|$ chute de 0.125 à 0.044. Ce gain conjoint en performance et en équité est une exception au compromis habituel : les proxies retirés portaient un signal corrélé au biais sans apporter de valeur prédictive substantielle, et leur suppression revient à diminuer la sur-paramétrisation du modèle. Sur un autre dataset, ce résultat pourrait ne pas se reproduire et il faudrait alors arbitrer explicitement entre performance et équité.
\end{leconbox}

% ════════════════════════════════════════════════════════════════
\section{Quantification d'incertitude}

\subsection{Bootstrap}

Suite au conseil de l'enseignante, le volet robustesse adversariale du sujet est remplacé par une quantification d'incertitude statistique via bootstrap. On entraîne $B = 500$ régressions logistiques sur des ré-échantillonnages avec remise du train (Efron 1979), puis on agrège les métriques par quantiles. Cette approche ne couvre pas toute l'incertitude de généralisation (le test reste fixe) mais elle quantifie précisément la sensibilité des conclusions au tirage de l'échantillon d'entraînement.

\begin{table}[H]
\centering
\begin{tabular}{lcccc}
\toprule
\textbf{Métrique} & \textbf{Moyenne} & $\sigma$ & \textbf{IC 2.5\,\%} & \textbf{IC 97.5\,\%} \\
\midrule
AUC                       & 0.764 & 0.020 & 0.724 & 0.798 \\
BalAcc                    & 0.681 & 0.023 & 0.636 & 0.724 \\
$|\Delta\text{DP}|$ genre & 0.035 & 0.027 & 0.001 & 0.104 \\
$|\Delta\text{EO}|$ genre & 0.089 & 0.042 & 0.020 & 0.177 \\
DI genre                  & 0.92  & 0.060 & 0.78  & 1.00  \\
$|\Delta\text{DP}|$ âge   & 0.103 & 0.050 & 0.008 & 0.205 \\
$|\Delta\text{EO}|$ âge   & 0.087 & 0.050 & 0.021 & 0.195 \\
DI âge                    & 0.77  & 0.109 & 0.56  & 0.98  \\
\bottomrule
\end{tabular}
\caption{Bootstrap $B = 500$ sur le baseline V2 (sans attributs sensibles).}
\end{table}

\begin{figure}[H]
\centering
\includegraphics[width=0.95\linewidth]{../outputs/05_bootstrap.png}
\end{figure}

\begin{attentionbox}[{Lecture critique}]
L'IC 95\,\% de $|\Delta\text{DP}|$ âge va de \textbf{0.01 à 0.21}. Le biais empirique (0.125) est dans cet IC, mais l'incertitude est plus large que la valeur. Annoncer \og{} biais $|\Delta\text{DP}| = 0.125$ \fg{} sans préciser \og{} $\pm 0.10$ \fg{} serait trompeur.
\end{attentionbox}

\subsection{Intervalle de prédiction sur un client unique}

\begin{figure}[H]
\centering
\includegraphics[width=0.95\linewidth]{../outputs/06_indiv_uncertainty.png}
\caption{À gauche : intervalle de confiance bootstrap par individu, trié par score moyen. À droite : écart-type bootstrap en fonction du score moyen, coloré par la part des modèles qui prédisent défaut.}
\end{figure}

L'analyse individuelle prolonge le bootstrap global au niveau du dossier : pour chaque client, on calcule la moyenne, l'écart-type et l'IC 95 \% des scores prédits par les $B$ modèles, ainsi que la part des modèles qui classent ce client en défaut ($p_{\text{décision défaut}}$). Le panneau de gauche montre la forme attendue (l'incertitude est maximale au voisinage du seuil de décision et diminue aux extrêmes), et le panneau de droite met en évidence la parabole d'incertitude caractéristique. Le notebook isole ensuite les 10 dossiers \emph{sensibles}, définis comme ceux dont l'IC bootstrap traverse le seuil $\tau$ ; pour ces clients, une décision purement automatique est instable et une revue humaine est recommandée.

\subsection{Bootstrap sur les 3 variantes}

\begin{table}[H]
\centering
\begin{tabular}{llccc}
\toprule
\textbf{Variante} & \textbf{Métrique} & \textbf{Moyenne} & \textbf{IC 2.5\,\%} & \textbf{IC 97.5\,\%} \\
\midrule
V1 (avec sensibles) & AUC                  & 0.764 & 0.721 & 0.796 \\
V1 (avec sensibles) & $|\Delta\text{DP}|$ âge & 0.126 & 0.024 & 0.230 \\
V1 (avec sensibles) & $|\Delta\text{EO}|$ âge & 0.111 & 0.026 & 0.239 \\
\midrule
V2 (sans sensibles) & AUC                  & 0.765 & 0.727 & 0.795 \\
V2 (sans sensibles) & $|\Delta\text{DP}|$ âge & 0.102 & 0.013 & 0.199 \\
V2 (sans sensibles) & $|\Delta\text{EO}|$ âge & 0.086 & 0.022 & 0.188 \\
\midrule
\textbf{V3 (sans + proxies)} & AUC                  & 0.768 & 0.728 & 0.799 \\
\textbf{V3 (sans + proxies)} & $|\Delta\text{DP}|$ âge & \textbf{0.046} & \textbf{0.004} & \textbf{0.127} \\
\textbf{V3 (sans + proxies)} & $|\Delta\text{EO}|$ âge & 0.086 & 0.025 & 0.175 \\
\bottomrule
\end{tabular}
\caption{Bootstrap $B = 100$ sur les 3 variantes. Cible : âge.}
\end{table}

\begin{figure}[H]
\centering
\includegraphics[width=0.95\linewidth]{../outputs/19_bootstrap_triple.png}
\caption{Distributions bootstrap des 3 variantes. AUC : largement chevauchantes (pas de coût significatif). $|\Delta\text{DP}|$ âge : V3 nettement plus bas et plus serré.}
\end{figure}

\begin{leconbox}[V3 plus stable que V2 et V1]
Sur l'AUC, les 3 distributions se recouvrent : retirer les attributs sensibles puis les proxies n'a pas de coût statistiquement significatif. Sur $|\Delta\text{DP}|$ âge, V3 a l'IC 95\,\% le plus serré et le plus bas. \textbf{V3 est préférable à la fois en moyenne (biais plus faible) et en stabilité statistique (IC plus serré).}
\end{leconbox}

% ════════════════════════════════════════════════════════════════
\section{Synthèse fairness et conclusion}

\begin{figure}[H]
\centering
\includegraphics[width=0.95\linewidth]{../outputs/12_dashboard.png}
\caption{Tableau de bord récapitulatif du baseline V2. Ligne pointillée rouge : seuil DI 80~\% (EEOC).}
\end{figure}

\subsection{Reproductibilité méthodologique}

Le but de cette section est de montrer que le canevas suivi est transférable à n'importe quel projet de classification binaire avec attributs sensibles. La démarche est conçue pour qu'un praticien puisse la réappliquer à un cas d'usage différent (scoring assurance, recrutement, allocation d'aides sociales, justice prédictive, scoring santé) en remplaçant essentiellement le dataset et la définition des attributs sensibles, sans avoir à réinventer la structure de l'audit.

Le canevas en 8 étapes : (i) identification des attributs sensibles et préparation des splits stratifiés ; (ii) audit du baseline en performance, équité par attribut et intersectionnel ; (iii) détection des proxies par corrélation univariée et régression multivariée features $\to$ attribut sensible ; (iv) comparaison de trois variantes V1 (avec sensibles), V2 (sans), V3 (sans + proxies) ; (v) évaluation des méthodes d'atténuation classiques (pré-traitement par reweighing, post-traitement par seuils par groupe) ; (vi) interprétabilité globale et locale ; (vii) quantification d'incertitude par bootstrap au niveau global et individuel ; (viii) synthèse et parcours décisionnel motivé.

Trois catégories de choix interviennent dans ce canevas et il faut les distinguer pour le transférer correctement. Les \emph{choix universels} gardent la même forme d'un projet à l'autre et constituent le squelette de la démarche : la structure en 8 étapes, le triptyque V1/V2/V3, l'utilisation conjointe de $|\Delta\text{DP}|$, $|\Delta\text{EO}|$ et DI, l'audit intersectionnel, le diagnostic des proxies par AUC features $\to$ attribut sensible, et le bootstrap pour quantifier l'incertitude. Les \emph{choix dépendants du domaine} doivent être révisés à chaque nouvelle application : la définition opérationnelle des attributs sensibles (ici, seuil 25 ans pour l'âge, suivant Kamiran \& Calders 2012) ; la matrice de coûts métier (ici, $5 \cdot \text{FN} + 1 \cdot \text{FP}$, documentée par UCI) ; le seuil légal de référence (ici, EEOC 80~\%, mais d'autres juridictions retiennent d'autres seuils) ; le critère d'équité prioritaire (DP pour la conformité légale, EO pour la justice individuelle, le choix dépend du contexte réglementaire et opérationnel). Les \emph{choix dépendants du dataset} doivent être recalibrés à chaque jeu particulier : le seuil de corrélation pour la détection de proxies ($|r| > 0.15$ ici, à ajuster selon la dimensionnalité et le bruit) ; la taille $B$ du bootstrap (500 ici, à augmenter sur des datasets plus grands ou plus rapides à entraîner) ; la famille du modèle baseline (régression logistique ici parce que Linear SHAP est exact, à remplacer par un Random Forest ou un Gradient Boosting si la non-linéarité est attendue, avec TreeSHAP en remplacement).

La conclusion principale du projet --- V3 améliore simultanément performance et équité par attribut isolé sans intervention in/post-modèle, mais ne résout pas le résidu intersectionnel --- est elle aussi spécifique à ce dataset. Sur un autre cas d'usage, l'arbitrage pourrait basculer : V3 pourrait coûter beaucoup en performance si les proxies portent du signal prédictif utile, ou ne rien améliorer en équité si le biais est de nature non-proxy (par exemple porté par la définition même de la cible). Ce qui est transférable, c'est le fait de toujours comparer V1, V2 et V3 explicitement plutôt que de supposer \emph{a priori} l'effet d'un retrait, et de quantifier le résultat avec ses intervalles de confiance avant de conclure.

\subsection{Bilan des résultats}

\paragraph{Performance.} Le baseline V2 (régression logistique L2) atteint sur le test AUC = 0.785, balanced accuracy = 0.669 et average précision = 0.607, avec une convergence en 6 itérations. La matrice de coûts asymétrique documentée par UCI ($5 \cdot \text{FN} + 1 \cdot \text{FP}$) ancre le choix du seuil de classification dans la réalité bancaire.

\paragraph{Équité par attribut isolé.} Sur le genre, le biais est quasi nul après retrait de \texttt{personal\_status\_sex} ($|\Delta\text{DP}| = 0.010$, DI = 0.97). Sur l'âge, en revanche, le biais persiste : $|\Delta\text{DP}| = 0.125$ et DI = 0.69, soit sous le seuil légal US 0.8. L'audit intersectionnel multiplie ce biais par environ 1.5 ($|\Delta\text{DP}| = 0.181$, DI = 0.59), avec les jeunes hommes (n = 19) comme sous-groupe le plus défavorisé, contre-intuitivement.

\paragraph{Comparaison des trois variantes.} La performance est strictement préservée par les retraits successifs (V1 = 0.781, V2 = 0.785, V3 = 0.790). Sur l'âge, $|\Delta\text{DP}|$ chute de 0.138 (V1) à 0.125 (V2) puis à 0.044 (V3), et DI passe de 0.62 à 0.69 puis à 0.87, au-dessus du seuil EEOC pour la première fois. Sur l'intersection genre $\times$ âge, le résultat est moins net : V3 ne fait pas mieux que V1 (DI = 0.48 contre 0.44), ce qui indique que l'équité intersectionnelle ne résulte pas mécaniquement de l'équité par attribut isolé.

\paragraph{Atténuation classique.} Le reweighing a un effet limité sur V1 et V2, et nul sur V3 dont les marges sont déjà équilibrées. Les seuils par groupe (post-processing) réduisent $|\Delta\text{DP}|$ sur V1 et V2 au prix d'une hausse de $|\Delta\text{EO}|$ (théorème d'impossibilité). Sur V3, ce même post-processing devient contre-productif : il fait remonter $|\Delta\text{DP}|$ de 0.044 à 0.124, parce qu'il force le modèle à s'écarter d'un équilibre déjà atteint pour viser une cible (égalité stricte des taux de sélection) que le baseline ne nécessitait plus.

\paragraph{Atténuation par retrait des proxies.} L'AUC de la régression features $\to$ âge atteint 0.86, ce qui démontre que retirer seulement \texttt{age\_in\_years} laisse subsister un signal proxy multivarié majeur. Le retrait additionnel des 17 features avec $|r| > 0.15$ produit V3, qui constitue le meilleur compromis global : meilleure AUC, meilleure équité par attribut isolé, et plus grande stabilité statistique, sans intervention in-modèle ni post-modèle.

\paragraph{Interprétabilité.} SHAP global et permutation importance concordent sur les trois variantes en désignant les mêmes variables financières dominantes (\texttt{checking\_status}, \texttt{credit\_history}, \texttt{installment\_rate}, \texttt{savings\_account\_bonds}). SHAP local décompose chaque décision en contributions par feature et fournit le matériau d'une motivation conforme à l'article 22 du RGPD.

\paragraph{Quantification d'incertitude.} Le bootstrap $B = 500$ sur V2 donne un IC 95~\% sur l'AUC de $[0.72, 0.80]$ et sur $|\Delta\text{DP}_{\text{age}}|$ de $[0.01, 0.21]$ : l'incertitude est plus large que la valeur ponctuelle, ce qui doit modérer toute conclusion. Le bootstrap $B = 100$ sur les trois variantes confirme que V3 est préférable à la fois en moyenne (biais central plus bas) et en stabilité (IC le plus serré, $[0.004, 0.127]$).

\subsection{Parcours décisionnel}

L'enchaînement des méthodes appliqué dans ce projet peut se résumer comme suit. On commence par auditer le baseline V2 et on détecte une alerte sur l'âge ($|\Delta\text{DP}_{\text{age}}| = 0.125$, DI = 0.69, sous le seuil légal). L'audit intersectionnel localise le biais sur le sous-groupe \emph{male\_younger} (n = 19), et la détection de proxies montre que retirer les attributs sensibles ne suffira pas (AUC features $\to$ âge = 0.86, 17 proxies univariés actifs). On construit alors V3 en retirant ces proxies, et la comparaison V1 vs V2 vs V3 confirme que V3 réduit $|\Delta\text{DP}_{\text{age}}|$ à 0.044 sans coût de performance. Les méthodes d'atténuation classiques (reweighing et post-processing) sont testées sur V3 par souci d'exhaustivité, mais elles s'avèrent inutiles ou contre-productives, parce que V3 est déjà à l'équilibre. Un dernier audit intersectionnel sur V3 met en évidence un résidu non résolu (DI = 0.48), qui appelle une décision explicite : accepter ce résidu, ou ajouter une mitigation in-modèle spécifique aux sous-groupes (non implémentée ici). Le bootstrap confirme finalement la robustesse statistique de la conclusion, et SHAP local fournit l'outil de motivation des décisions individuelles requis par le RGPD.

\subsection{Limites}

L'étude porte sur un dataset historiquement utilisé pour ce type d'évaluation mais de taille réduite (1000 lignes, test 200), ce qui produit des IC 95~\% sur l'équité dépassant souvent 0.15 ; cette incertitude doit être communiquée explicitement avec toute conclusion. L'approche reste linéaire (régression logistique), avec une capacité limitée à capturer des interactions non linéaires entre features. Le reweighing utilisé ne traite qu'une dépendance binaire $S \perp Y$ et serait à généraliser pour un cadre multi-attribut. Le retrait massif de 17 proxies en V3 est une stratégie agressive : sur un autre dataset, ces features pourraient porter du signal prédictif utile et leur suppression dégraderait la performance. L'intersection genre $\times$ âge reste non résolue par V3, ce qui montre les limites d'une approche fondée uniquement sur l'équité par attribut. Enfin, la Fairness through Awareness (Dwork et al.\ 2012), qui aurait été pertinente compte tenu de l'AUC proxy = 0.86, n'a pas été implémentée.

\subsection{Recommandations pour un déploiement réel}

Pour un usage opérationnel sur ce dataset, nous recommandons d'adopter V3 (sans attributs sensibles, sans les 17 proxies univariés $|r| > 0.15$) comme baseline, car elle constitue le meilleur compromis performance/équité par attribut et présente la plus grande stabilité statistique. Le post-processing correctif ne doit pas être appliqué sur V3, sous peine de dégrader l'équité. Le résidu intersectionnel (DI $\approx 0.48$) doit être documenté et accompagné d'une procédure de revue manuelle pour les sous-groupes minoritaires. Le critère d'équité prioritaire (DP pour la conformité EEOC, EO pour la justice individuelle) doit être choisi explicitement et fixé contractuellement avant déploiement. Les dossiers pour lesquels l'IC bootstrap traverse le seuil de décision (\texttt{traverse\_seuil = True} dans le notebook) doivent systématiquement déclencher une expertise humaine. L'analyse intersectionnelle gagnerait à être reproduite sur un dataset plus grand pour confirmer les conclusions sur les sous-groupes minoritaires.

% ════════════════════════════════════════════════════════════════
\section*{Références}

{\footnotesize
\begin{thebibliography}{99}
\itemsep=0.1em
\bibitem{chouldechova2017} A.~Chouldechova, \textit{Fair prediction with disparate impact}, Big Data 5(2), 2017. \href{https://arxiv.org/abs/1703.00056}{arXiv:1703.00056}.
\bibitem{kleinberg2017} J.~Kleinberg, S.~Mullainathan, M.~Raghavan, \textit{Inherent trade-offs in the fair determination of risk scores}, ITCS 2017. \href{https://arxiv.org/abs/1609.05807}{arXiv:1609.05807}.
\bibitem{hardt2016} M.~Hardt, E.~Price, N.~Srebro, \textit{Equality of Opportunity in Supervised Learning}, NeurIPS 2016. \href{https://arxiv.org/abs/1610.02413}{arXiv:1610.02413}.
\bibitem{kamiran2012} F.~Kamiran, T.~Calders, \textit{Data preprocessing techniques for classification without discrimination}, KAIS 33(1), 2012. \href{https://doi.org/10.1007/s10115-011-0463-8}{doi}.
\bibitem{dwork2012} C.~Dwork et al., \textit{Fairness through awareness}, ITCS 2012. \href{https://arxiv.org/abs/1104.3913}{arXiv:1104.3913}.
\bibitem{bellamy2018} R.~K.~E.~Bellamy et al., \textit{AI Fairness 360}, IBM, 2018. \href{https://arxiv.org/abs/1810.01943}{arXiv:1810.01943}.
\bibitem{bird2020} S.~Bird et al., \textit{Fairlearn}, Microsoft 2020. \href{https://fairlearn.org/}{fairlearn.org}.
\bibitem{lundberg2017} S.~Lundberg, S.-I.~Lee, \textit{A Unified Approach to Interpreting Model Predictions}, NeurIPS 2017. \href{https://arxiv.org/abs/1705.07874}{arXiv:1705.07874}. Linear SHAP : Corollary 1.
\bibitem{efron1979} B.~Efron, \textit{Bootstrap methods}, Annals of Statistics 7(1), 1979.
\bibitem{lakshminarayanan2017} B.~Lakshminarayanan et al., \textit{Deep Ensembles}, NeurIPS 2017. \href{https://arxiv.org/abs/1612.01474}{arXiv:1612.01474}.
\bibitem{pedregosa2011} F.~Pedregosa et al., \textit{Scikit-learn}, JMLR 12, 2011.
\bibitem{cours} Cours IA708, F.~d'Alché-Buc, C.~Laclau, G.~Franchi, Télécom Paris.
\end{thebibliography}
}

\end{document}
